\chapter{The oscilloscope}

\section{One lane per voice}
The oscilloscope draws the shape of each voice's sound. Each chip takes a column, each voice a row, with its name at
the top of the cell. You see who plays what: the bass that holds, the arpeggio that shimmers, the snare that flashes
past.

To keep shapes still, we lock each lane's window to the period of its note: this is triggering, as on a lab
oscilloscope. We then reduce the signal to two points per pixel column, the minimum and the maximum, so no peak is
lost even when the lane is narrow.

\section{Special cases}
With the SID samples lane on, a fourth row per chip shows the samples played through the volume. For the AdLib, we lay
out the nine melodic voices 3 by 3 and hide the five rhythm-mode lanes, which few tracks use; turn on ``AdLib rhythm
lanes'' if you need them.

\section{SID filter response}
The visualisation button replaces the top half of the oscilloscope with the filter curve of each SID, with each
voice's position. Previous states stay visible for a moment as they fade, which makes fast cutoff sweeps readable.
\capturedouble{ux4-lowpass-waves}{ux4-sweep}{Resonant low-pass filter}{cutoff sweep}
If you want more room for the voices, reduce the filter response share in the oscilloscope settings.
\capture[0.7]{ux4-lowpass-share35}{Filter response reduced to 35~\% of the height}

\section{Settings}
\capture{prefs-en-oscilloscope}{Oscilloscope settings}
\begin{reglages}
\regl{Waveform window}{5 to 100~ms, \textbf{40~ms}}{Length of sound shown per voice. 40~ms is the value of corrscope
(\texttt{render\_ms}), the tool behind most chiptune oscilloscope videos.}
\regl{Visual gain}{\textbf{automatic per voice}, $\times$1 to $\times$8}{Affects only the drawing. In automatic mode
each lane is scaled to its own level, so a quiet voice stays readable.}
\regl{Maximum periods shown}{off, 4, \textbf{8}, 16, 32}{Zooms in on high notes, which otherwise shrink to a thick
line. Specific to \soft.}
\regl{SID gauges window}{2 to 40~ms, \textbf{10~ms}}{Length summarised by the SID oscillator gauges.}
\regl{Oscilloscope height}{60 to 160~\%, \textbf{100~\%}}{Relative to a 16:9 shape of the width; the dashboard takes
the rest.}
\regl{Filter response height}{20 to 80~\%, \textbf{55~\%}}{Share of the SID filter view taken by the curve.}
\regl{Waveforms under the filter}{no, \textbf{yes}}{Otherwise the curve alone.}
\regl{AdLib rhythm lanes}{\textbf{no}, yes}{See above. Applied to the next track.}
\regl{Frame rate}{\textbf{screen rate}, 120, 90, 60, 30}{A lower limit saves battery.}
\regl{Graphics engine}{\textbf{automatic}, Vulkan, OpenGL ES}{Automatic picks Vulkan when available.}
\regl{Custom Vulkan driver}{\textbf{no}, package to choose}{On Adreno GPUs, a Mesa Turnip driver can replace the
vendor's, as recent Android emulators do. The system driver takes over if yours fails to load.}
\end{reglages}

\section{Triggering}
Locking a waveform so it stays still is harder than it looks: a note changing timbre, an arpeggio, noise have no clean
edge. Lab oscilloscopes look for a threshold crossing; that is not enough for music. The state of the art is
corrscope's correlation trigger, by nyanpasu64: it looks, in the incoming sound, for the position most similar to the
previous frames, favouring rising edges. We ported it as is, with its default values, so that people who know
corrscope get the same picture and can reuse their settings.
\begin{reglages}
\regl{Preset}{\textbf{corrscope default (40~ms)}, evolving notes (60~ms), chiptune square, basses, NES triangle,
power saving, custom}{Each preset sets the settings below.}
\regl{Width}{10 to 200~ms, \textbf{40~ms}}{corrscope's \texttt{trigger\_ms}. Must cover two periods of the lowest
notes: raise it for basses.}
\regl{Edge strength}{0 to 300~\%, \textbf{100~\%}}{Favours steep rising edges: new notes are caught better, held
notes move more.}
\regl{Slope width}{5 to 100~\% of the period, \textbf{25~\%}}{Width of the edge looked for.}
\regl{Buffer strength}{0 to 300~\%, \textbf{100~\%}}{Favours similarity with previous frames.}
\regl{Responsiveness}{0 to 100~\%, \textbf{50~\%}}{How fast the history follows the sound. corrscope expresses it per
frame at 60~fps; we convert it for other screen rates.}
\regl{Falloff}{10 to 200~\% of the period, \textbf{50~\%}}{Weight of the history's edges.}
\regl{Reset below}{0 to 100~\%, \textbf{30~\%}}{Forgets the history when the new note no longer resembles it.}
\regl{DC removal}{0 to 100~\%, \textbf{0}}{For lanes whose centre drifts.}
\regl{Sign amplification}{0 to 300~\%, \textbf{0}}{Creates sharp zero crossings for soft waves such as the triangle.}
\regl{Search range}{10 to 100~\% of the width, \textbf{50~\%}}{How far the shape may jump from one frame to the
next.}
\regl{Falling edges}{\textbf{no}, yes}{For sounds that start with a fall, like some sampled trumpets.}
\regl{Pitch tracking}{no, \textbf{yes}}{Rescales the history when the note changes.}
\regl{Zero-crossing snap}{\textbf{0} samples}{Slides to the nearest rising zero crossing.}
\regl{Full searches}{\textbf{every frame}, 60, 30 per second}{Fewer searches, less computation.}
\regl{Subsampling}{\textbf{automatic}, none, 2, 4}{Computes the correlation on fewer points.}
\end{reglages}
